% questions.tex —— 2014 年·重庆卷·理科 · 题目内容片段（无答案）
% 由 tools/generate_projects.py 依来源总账自 2014/2014_数学真题_重庆理_试题.md 生成，勿手改（改源 MD 后重跑）


\begin{gknotice}
\textbf{来源说明}：本题面依据固定源 PDF 的两页 A3 题面逐题转录，原卷印刷页码为 865—866。第 865 面为选择题 1—10、填空题 11—16 与解答题 17—18；第 866 面为解答题 19—22。全卷三处节头依次为「一、选择题」「二、填空题」「三、解答题」；题号与索引一致。

\textbf{解析归属}：固定源不附答案或解析；配套答案与解析均由本库按本题面独立复算。

\textbf{图形说明}：第 5 题的程序框图、第 7 题的三视图、第 19 题的四棱锥、第 21 题的椭圆与焦点图由项目内 TikZ 矢量源重绘；其余各题原页无图。
\end{gknotice}

\gkpart{一、选择题}

\begin{question}
\gkti{1}在复平面内表示复数 $i(1-2i)$ 的点位于\mbox{（　　）}

（A）第一象限；（B）第二象限；（C）第三象限；（D）第四象限。

\end{question}
\begin{question}
\gkti{2}对任意等比数列 $\{a_n\}$，下列说法一定正确的是\mbox{（　　）}

（A）$a_1$，$a_3$，$a_9$ 成等比数列；（B）$a_2$，$a_3$，$a_6$ 成等比数列；（C）$a_2$，$a_4$，$a_8$ 成等比数列；（D）$a_3$，$a_6$，$a_9$ 成等比数列。

\end{question}
\begin{question}
\gkti{3}已知变量 $x$ 与 $y$ 正相关，且由观测数据算得样本平均数 $\bar x=3$，$\bar y=3.5$，则由该观测数据算得的线性回归方程可能是\mbox{（　　）}

（A）$\hat y=0.4x+2.3$；（B）$\hat y=2x-2.4$；（C）$\hat y=-2x+9.5$；（D）$\hat y=-0.3x+4.4$。

\end{question}
\begin{question}
\gkti{4}已知向量 $\boldsymbol a=(k,3)$，$\boldsymbol b=(1,4)$，$\boldsymbol c=(2,1)$，且 $(2\boldsymbol a-3\boldsymbol b)\perp\boldsymbol c$，则实数 $k=$\mbox{（　　）}

（A）$-\dfrac92$；（B）$0$；（C）$3$；（D）$\dfrac{15}2$。

\Needspace{0.4\textheight}
\end{question}
\begin{question}
\gkti{5}执行如图所示的程序框图，若输出 $k$ 的值为 $6$，则判断框内可填入的条件是\mbox{（　　）}

（A）$s>\dfrac12$；（B）$s>\dfrac35$；（C）$s>\dfrac7{10}$；（D）$s>\dfrac45$。
\nopagebreak[4]
\begin{center}% ============================================================
% f2014_chongqing_science_q05_flowchart.tex
% 2014 重庆卷（理科）第 5 题：程序框图
% 来源 PDF：普通高考/2014/2014重庆理.pdf
% 原页：第 1 页（A3 横排）左栏，印面 865，第 5 题题干下方
% 结构：开始 → 赋初值 k=9、s=1 → 判断框（框内待填条件）；
%       是：s=s·k/(k+1) → k=k-1 → 返回判断；否：输出 k → 结束。
% ============================================================
\begin{tikzpicture}[>=Stealth, line width=0.7pt, font=\small,
  line cap=round, line join=round]

  % 开始
  \draw[rounded corners=6pt] (-0.95,0.45) rectangle (0.95,-0.45);
  \node at (0,0) {开始};
  \draw[->] (0,-0.45) -- (0,-0.95);

  % 初始化
  \draw (-1.75,-0.95) rectangle (1.75,-1.85);
  \node at (0,-1.40) {$k=9$，$s=1$};

  % 进入判断框
  \draw[->] (0,-1.85) -- (0,-2.85);

  % 判断菱形（原页框内为空，条件待填）
  \draw (0,-2.85) -- (2.15,-3.90) -- (0,-4.95) -- (-2.15,-3.90) -- cycle;

  % 否：向下输出
  \draw[->] (0,-4.95) -- (0,-5.45);
  \node at (-0.38,-5.20) {否};
  \draw (-1.95,-5.45) -- (1.55,-5.45) -- (1.95,-6.45) -- (-1.55,-6.45) -- cycle;
  \node at (0,-5.95) {输出 $k$};
  \draw[->] (0,-6.45) -- (0,-6.95);

  % 结束
  \draw[rounded corners=6pt] (-0.95,-6.95) rectangle (0.95,-7.85);
  \node at (0,-7.40) {结束};

  % 是：向右进入循环体
  \draw[->] (2.15,-3.90) -- (2.95,-3.90);
  \node at (2.55,-3.58) {是};

  % 循环体第一框
  \draw (2.95,-4.40) rectangle (6.45,-3.40);
  \node at (4.70,-3.90) {$s=s\cdot\dfrac{k}{k+1}$};

  % 循环体第二框
  \draw[->] (4.70,-3.40) -- (4.70,-2.70);
  \draw (2.95,-2.70) rectangle (6.45,-1.70);
  \node at (4.70,-2.20) {$k=k-1$};

  % 回边：向左汇入进入判断框的竖线
  \draw (2.95,-2.20) -- (0,-2.20);
\end{tikzpicture}
\end{center}
\nopagebreak[4]
\end{question}
\begin{question}
\gkti{6}已知命题 $p$：对任意 $x\in\mathbf R$，总有 $2^x>0$；$q$：“$x>1$”是“$x>2$”的充分不必要条件。则下列命题为真命题的是\mbox{（　　）}

（A）$p\wedge q$；（B）$\neg p\wedge\neg q$；（C）$\neg p\wedge q$；（D）$p\wedge\neg q$。

\Needspace{0.4\textheight}
\end{question}
\begin{question}
\gkti{7}某几何体的三视图如图所示，则该几何体的表面积为\mbox{（　　）}

（A）$54$；（B）$60$；（C）$66$；（D）$72$。
\nopagebreak[4]
\begin{center}% ============================================================
% f2014_chongqing_science_q07_three_views.tex
% 2014 重庆卷（理科）第 7 题：三视图
% 来源 PDF：普通高考/2014/2014重庆理.pdf
% 原页：第 1 页（A3 横排）上栏，印面 865，第 7 题题干上方
% 说明：正视图为左高 5、右高 2、底长 4 的直角梯形；左视图为
%       宽 3、高 5 的矩形，内含一条从左上角斜向右下的虚线；
%       俯视图为两直角边长 4 与 3 的直角三角形。
% ============================================================
\begin{tikzpicture}[>=Stealth, line width=0.7pt, font=\small,
  x=1cm, y=1cm, line cap=round, line join=round,
  ext/.style={line width=0.4pt}]

  % ---------- 正视图 ----------
  \draw (0,0) -- (4,0) -- (4,2) -- (0,5) -- cycle;
  % 高 5
  \draw[ext] (0,0) -- (-0.6,0);
  \draw[ext] (0,5) -- (-0.6,5);
  \draw[<->] (-0.6,0) -- (-0.6,5) node[midway,left,inner sep=1pt] {$5$};
  % 底 4
  \draw[ext] (0,0) -- (0,-0.6);
  \draw[ext] (4,0) -- (4,-0.6);
  \draw[<->] (0,-0.6) -- (4,-0.6) node[midway,below,inner sep=2pt] {$4$};
  % 右高 2
  \draw[ext] (4,0) -- (4.6,0);
  \draw[ext] (4,2) -- (4.6,2);
  \draw[<->] (4.6,0) -- (4.6,2) node[midway,right,inner sep=1pt] {$2$};
  \node at (2,-1.5) {正视图};

  % ---------- 左视图 ----------
  \draw (6,0) rectangle (9,5);
  \draw[dashed] (6,5) -- (9,2);
  \draw[ext] (6,0) -- (6,-0.6);
  \draw[ext] (9,0) -- (9,-0.6);
  \draw[<->] (6,-0.6) -- (9,-0.6) node[midway,below,inner sep=2pt] {$3$};
  \node at (7.5,-1.5) {左视图};

  % ---------- 俯视图 ----------
  \draw (11,0) -- (15,0) -- (11,3) -- cycle;
  \node at (13,-1.5) {俯视图};
\end{tikzpicture}
\end{center}
\nopagebreak[4]
\end{question}
\begin{question}
\gkti{8}设 $F_1$，$F_2$ 分别为双曲线 $\dfrac{x^2}{a^2}-\dfrac{y^2}{b^2}=1$（$a>0$，$b>0$）的左、右焦点，双曲线上存在一点 $P$ 使得 $|PF_1|+|PF_2|=3b$，$|PF_1|\cdot|PF_2|=\dfrac94 ab$，则该双曲线的离心率为\mbox{（　　）}

（A）$\dfrac43$；（B）$\dfrac53$；（C）$\dfrac94$；（D）$3$。

\end{question}
\begin{question}
\gkti{9}某次联欢会要安排 $3$ 个歌舞类节目、$2$ 个小品类节目和 $1$ 个相声类节目的演出顺序，则同类节目不相邻的排法种数是\mbox{（　　）}

（A）$72$；（B）$120$；（C）$144$；（D）$168$。

\end{question}
\begin{question}
\gkti{10}已知 $\triangle ABC$ 的内角 $A$，$B$，$C$ 满足 $\sin 2A+\sin(A-B+C)=\sin(C-A-B)+\dfrac12$，面积 $S$ 满足 $1\leqslant S\leqslant 2$，记 $a$，$b$，$c$ 分别为 $A$，$B$，$C$ 所对的边，则下列不等式一定成立的是\mbox{（　　）}

（A）$bc(b+c)>8$；（B）$ab(a+b)>16\sqrt2$；（C）$6\leqslant abc\leqslant 12$；（D）$12\leqslant abc\leqslant 24$。

\end{question}
\gkpart{二、填空题}

\begin{question}
\gkti{11}设全集 $U=\{n\in\mathbf N\mid 1\leqslant n\leqslant 10\}$，$A=\{1,2,3,5,8\}$，$B=\{1,3,5,7,9\}$，则 $(\complement_U A)\cap B=$\gkblank{}。

\end{question}
\begin{question}
\gkti{12}函数 $f(x)=\log_2\sqrt x\cdot\log_{\sqrt2}(2x)$ 的最小值为\gkblank{}。

\end{question}
\begin{question}
\gkti{13}已知直线 $ax+y-2=0$ 与圆心为 $C$ 的圆 $(x-1)^2+(y-a)^2=4$ 相交于 $A$，$B$ 两点，且 $\triangle ABC$ 为等边三角形，则实数 $a=$\gkblank{}。

\end{question}
\begin{question}
\gkti{14}过圆外一点 $P$ 作圆的切线 $PA$（$A$ 为切点），再作割线 $PBC$ 依次交圆于 $B$，$C$，若 $PA=6$，$AC=8$，$BC=9$，则 $AB=$\gkblank{}。

\end{question}
\begin{question}
\gkti{15}已知直线 $l$ 的参数方程为 $\begin{cases}x=2+t,\\y=3+t,\end{cases}$（$t$ 为参数），以坐标原点为极点，$x$ 轴的正半轴为极轴建立极坐标系，曲线 $C$ 的极坐标方程为 $\rho\sin^2\theta-4\cos\theta=0$（$\rho\geqslant 0$，$0\leqslant\theta<2\pi$），则直线 $l$ 与曲线 $C$ 的公共点的极径 $\rho=$\gkblank{}。

\end{question}
\begin{question}
\gkti{16}若不等式 $|2x-1|+|x+2|\geqslant a^2+\dfrac12 a+2$ 对任意实数 $x$ 恒成立，则实数 $a$ 的取值范围是\gkblank{}。

\end{question}
\gkpart{三、解答题}

\begin{question}
\gkti{17}已知函数 $f(x)=\sqrt3\sin(\omega x+\varphi)$（$\omega>0$，$-\dfrac\pi2\leqslant\varphi<\dfrac\pi2$）的图象关于直线 $x=\dfrac\pi3$ 对称，且图象上相邻两个最高点的距离为 $\pi$。

（1） 求 $\omega$ 和 $\varphi$ 的值；

（2） 若 $f\left(\dfrac\alpha2\right)=\dfrac{\sqrt3}4$（$\dfrac\pi6<\alpha<\dfrac{2\pi}3$），求 $\cos\left(\alpha+\dfrac{3\pi}2\right)$ 的值。

\end{question}
\begin{question}
\gkti{18}一盒中装有 $9$ 张各写有一个数字的卡片，其中 $4$ 张卡片上的数字是 $1$，$3$ 张卡片上的数字是 $2$，$2$ 张卡片上的数字是 $3$，从盒中任取 $3$ 张卡片。

（1） 求所取 $3$ 张卡片上的数字完全相同的概率；

（2） $X$ 表示所取 $3$ 张卡片上的数字的中位数，求 $X$ 的分布列与数学期望。

（注：若三个数 $a$，$b$，$c$ 满足 $a\leqslant b\leqslant c$，则称 $b$ 为这三个数的中位数）

\Needspace{0.4\textheight}
\end{question}
\begin{question}
\gkti{19}如图，四棱锥 $P-ABCD$ 中，底面是以 $O$ 为中心的菱形，$PO\perp$ 底面 $ABCD$，$AB=2$，$\angle BAD=\dfrac\pi3$，$M$ 为 $BC$ 上一点，且 $BM=\dfrac12$，$MP\perp AP$。
\nopagebreak[4]
\begin{center}% ============================================================
% f2014_chongqing_science_q19_pyramid.tex
% 2014 重庆卷（理科）第 19 题：四棱锥 P-ABCD
% 来源 PDF：普通高考/2014/2014重庆理.pdf
% 原页：第 1 页（A3 横排）右栏上部，印面 865，第 19 题题干右侧
% 说明：底面菱形 ABCD 画成平行四边形，O 为中心，顶点 P 在 O 的
%       正上方；M 为棱 BC 上靠近 B 的四等分点。被底面遮挡的
%       AD、DC 与对角线 AC、BD 以及 PD、PO 用虚线。
% ============================================================
\begin{tikzpicture}[>=Stealth, line width=0.7pt, font=\small,
  x=1cm, y=1cm, line cap=round, line join=round]

  \coordinate (A) at (0,2.2);
  \coordinate (B) at (2.6,0.4);
  \coordinate (C) at (7.6,1.4);
  \coordinate (D) at (5.0,3.2);
  \coordinate (O) at (3.8,1.8);
  \coordinate (P) at (3.7,5.8);
  \coordinate (M) at (3.85,0.65);

  % 被遮挡的棱
  \draw[dashed] (A) -- (D) -- (C);
  \draw[dashed] (A) -- (C);
  \draw[dashed] (B) -- (D);
  \draw[dashed] (P) -- (D);
  \draw[dashed] (P) -- (O);

  % 可见的棱
  \draw (A) -- (B) -- (C);
  \draw (P) -- (A);
  \draw (P) -- (B);
  \draw (P) -- (C);
  \draw (P) -- (M);

  \fill (O) circle (1.4pt);
  \node[below right, inner sep=3pt] at (O) {$O$};
  \node[left, inner sep=2pt] at (A) {$A$};
  \node[below, inner sep=3pt] at (B) {$B$};
  \node[right, inner sep=2pt] at (C) {$C$};
  \node[above left, inner sep=2pt] at (D) {$D$};
  \node[above, inner sep=3pt] at (P) {$P$};
  \node[below, inner sep=4pt] at (M) {$M$};
\end{tikzpicture}
\end{center}
\nopagebreak[4]
（1） 求 $PO$ 的长；

（2） 求二面角 $A-PM-C$ 的正弦值。

\end{question}
\begin{question}
\gkti{20}已知函数 $f(x)=ae^{2x}-be^{-2x}-cx$（$a$，$b$，$c\in\mathbf R$）的导函数 $f'(x)$ 为偶函数，且曲线 $y=f(x)$ 在点 $(0,f(0))$ 处的切线的斜率为 $4-c$。

（1） 确定 $a$，$b$ 的值；

（2） 若 $c=3$，判断 $f(x)$ 的单调性；

（3） 若 $f(x)$ 有极值，求 $c$ 的取值范围。

\Needspace{0.4\textheight}
\end{question}
\begin{question}
\gkti{21}设椭圆 $\dfrac{x^2}{a^2}+\dfrac{y^2}{b^2}=1$（$a>b>0$）的左、右焦点分别为 $F_1$，$F_2$，点 $D$ 在椭圆上，$DF_1\perp F_1F_2$，$\dfrac{|F_1F_2|}{|DF_1|}=2\sqrt2$，$\triangle DF_1F_2$ 的面积为 $\dfrac{\sqrt2}2$。
\nopagebreak[4]
\begin{center}% ============================================================
% f2014_chongqing_science_q21_ellipse.tex
% 2014 重庆卷（理科）第 21 题：椭圆与焦点
% 来源 PDF：普通高考/2014/2014重庆理.pdf
% 原页：第 2 页（A3 横排）中栏，印面 866，第 21 题题干下方
% 说明：椭圆焦点在 x 轴上，左、右焦点为 F1、F2；点 D 在椭圆上
%       且在 F1 的正下方，DF1⊥F1F2；图中画出线段 DF2。
% ============================================================
\begin{tikzpicture}[>=Stealth, line width=0.7pt, font=\small,
  x=1cm, y=1cm, line cap=round, line join=round]

  % 坐标轴
  \draw[->] (-5.4,0) -- (5.6,0) node[below, inner sep=3pt] {$x$};
  \draw[->] (0,-3.2) -- (0,3.4) node[right, inner sep=2pt] {$y$};

  % 椭圆
  \draw (0,0) ellipse (4 and 2.4);

  % 焦点与点 D
  \coordinate (F1) at (-2,0);
  \coordinate (F2) at (2,0);
  \coordinate (D) at (-2,-2.0785);

  \draw (F1) -- (D);
  \draw (D) -- (F2);

  \fill (F1) circle (1.4pt);
  \fill (F2) circle (1.4pt);
  \fill (D) circle (1.4pt);

  \node[above left, inner sep=2pt] at (F1) {$F_1$};
  \node[above right, inner sep=2pt] at (F2) {$F_2$};
  \node[below left, inner sep=2pt] at (D) {$D$};
  \node[below right, inner sep=2pt] at (0,0) {$O$};
\end{tikzpicture}
\end{center}
\nopagebreak[4]
（1） 求该椭圆的标准方程；

（2） 设圆心在 $y$ 轴上的圆与椭圆在 $x$ 轴的上方有两个交点，且圆在这两个交点处的两条切线相互垂直并分别过不同的焦点，求圆的半径。

\end{question}
\begin{question}
\gkti{22}设 $a_1=1$，$a_{n+1}=\sqrt{a_n^2-2a_n+2}+b$（$n\in\mathbf N^*$）。

（1） 若 $b=1$，求 $a_2$，$a_3$ 及数列 $\{a_n\}$ 的通项公式；

（2） 若 $b=-1$，问：是否存在实数 $c$，使得 $a_{2n}<c<a_{2n+1}$ 对所有 $n\in\mathbf N^*$ 成立？证明你的结论。

\end{question}